\documentclass[12pt]{article}

\usepackage{sbc-template}
\usepackage{graphicx,url}
\graphicspath{{imagens/}}
\usepackage[brazil]{babel}
\usepackage[utf8]{inputenc}
\usepackage{enumitem}
\usepackage{placeins}

\sloppy

\title{Avaliação de Desempenho e Sobrecarga de Protocolo:\\
       TCP, UDP e QUIC com NGINX e Linux NetEm}

\author{
  João Marcos Sousa Rufino Leal\inst{1}
}

\address{Departamento de Computação -- Universidade Federal do Piauí (UFPI)\\
  Bacharelado em Sistemas de Informação\\
  \email{jsousarufinoleal@ufpi.edu.br}
  \vspace{0.2cm}
  \newline
  \textbf{Vídeo de Demonstração e Apresentação (15 min):}\\
  \url{https://youtu.be/SEU_LINK_AQUI}
}

\begin{document}

\maketitle

\begin{abstract}
% A ser preenchido conforme o desenvolvimento do trabalho.
\end{abstract}

\begin{resumo}
% A ser preenchido conforme o desenvolvimento do trabalho.
\end{resumo}


% ============================================================
\section{Introdução e Objetivos}
% ============================================================
% A ser preenchido


% ============================================================
\section{Fundamentação Teórica}
% ============================================================

\subsection{TCP — RFC 793}
% A ser preenchido

\subsection{UDP Puro — RFC 768}
% A ser preenchido

\subsection{QUIC (RFC 9000/9001) e TLS 1.3 (RFC 8446)}
% A ser preenchido


% ============================================================
\section{Metodologia Experimental e Topologia}
% ============================================================
% A ser preenchido


% ============================================================
\section{Resultados Experimentais e Discussão}
% ============================================================

\subsection{Cenário A — Desempenho do UDP Puro em Baixa Latência}
% A ser preenchido

\subsection{Cenário B — Eficiência do TCP em Transferências Massivas}
% A ser preenchido

\subsection{Cenário C — Resiliência do QUIC em Redes Instáveis}
% A ser preenchido


% ============================================================
\section{Considerações Finais e Conclusão}
% ============================================================
% A ser preenchido


% ============================================================
\nocite{*}
\bibliographystyle{sbc}
\bibliography{sbc-template}
% ============================================================

\end{document}
